\documentclass[11pt]{article}
\usepackage[margin=1in]{geometry}
\usepackage[T1]{fontenc}
\usepackage{lmodern}
\usepackage{microtype}
\usepackage{amsmath,amssymb,amsthm,mathtools}
\usepackage{tikz-cd}

\setlength{\parindent}{0pt}
\setlength{\parskip}{0.45\baselineskip}
\allowdisplaybreaks

\newtheorem{theorem}{Theorem}[section]
\newtheorem{proposition}[theorem]{Proposition}
\newtheorem{corollary}[theorem]{Corollary}
\newtheorem{lemma}[theorem]{Lemma}
\theoremstyle{definition}
\newtheorem{remark}[theorem]{Remark}
\newtheorem{notation}[theorem]{Notation}

\newcommand{\Cat}{\mathrm{Cat}_{\infty}}
\newcommand{\C}{\mathcal{C}}
\newcommand{\Sps}{\mathcal{S}}
\newcommand{\PSh}{\mathcal{P}(\mathcal{C})}
\newcommand{\E}{\mathcal{E}}
\newcommand{\Ep}{\mathcal{E}^{\prime}}
\newcommand{\U}{\mathcal{U}}
\newcommand{\V}{\mathcal{V}}
\newcommand{\Fun}{\mathrm{Fun}}
\newcommand{\Map}{\mathrm{Map}}
\newcommand{\LFib}{\mathrm{LFib}}
\newcommand{\St}{\mathrm{St}}
\newcommand{\Un}{\mathrm{Un}}
\newcommand{\id}{\mathrm{id}}
\newcommand{\op}{\mathrm{op}}
\newcommand{\ev}{\mathrm{ev}}
\newcommand{\pt}{\ast}

\title{The $\infty$-Categorical Yoneda Lemma}
\author{}
\date{}

\begin{document}
\maketitle
\thispagestyle{empty}

\begin{center}
\emph{Note to the Reader}
\end{center}

The result proved below is a formal consequence of the straightening--
unstraightening equivalence, the universal property of sections of a left
fibration, and the finality of the terminal object of a slice $\infty$-category.
We work throughout in the quasicategorical model and assume familiarity with the
straightening formalism of \emph{Higher Topos Theory}, \S2.2.1, and the finality
criterion of \emph{Higher Topos Theory}, \S4.1.3.1.

\section{Conventions}

We work in the quasicategorical model of $\infty$-categories. Thus an
$\infty$-category is a simplicial set satisfying the inner horn-filling
condition. Equivalences are equivalences in the sense of $\infty$-categories;
limits are limits in the ambient $\infty$-categorical sense, hence homotopy
limits; mapping spaces are the canonical Kan complexes attached to an
$\infty$-category. Fix Grothendieck universes $\U \in \V$. The adjective
``small'' means ``$\U$-small.'' The $\infty$-category $\Sps$ of spaces is taken in
$\V$, with the full subcategory of $\U$-small spaces understood implicitly.
Consequently, if $\C$ is a small $\infty$-category, then
\[
  \PSh := \Fun(\C^{\op},\Sps)
\]
is a $\V$-small $\infty$-category.

For an $\infty$-category $K$, we write $\LFib(K)$ for the full subcategory of
$\mathrm{Cat}_{\infty,/K}$ spanned by left fibrations over $K$. If $p \colon X \to K$ is a
left fibration and $k \in K$ is an object, its fiber at $k$ means the pullback
\[
  X_k := X \times_K \{k\},
\]
regarded as a space. If $r \colon X \to K$ is any morphism in $\mathrm{Cat}_{\infty}$, we
write
\[
  \Gamma(r) := \Map_{\mathrm{Cat}_{\infty,/K}}(\id_K,r)
\]
for the space of sections of $r$.

We fix once and for all the straightening--unstraightening equivalence in the
form
\[
  \St_K \colon \LFib(K) \xrightarrow{\simeq} \Fun(K,\Sps),
  \qquad
  \Un_K \colon \Fun(K,\Sps) \xrightarrow{\simeq} \LFib(K),
\]
functorial in $K$; cf. HTT, \S2.2.1. All identifications below are taken with
respect to this equivalence and the universal properties it exhibits.

\section{Representable left fibrations}

Let $\C$ be a small $\infty$-category. For $C \in \C$, define the representable
presheaf
\[
  j(C) := \Map_{\C}(-,C) \in \Fun(\C^{\op},\Sps).
\]
This construction is functorial in $C$ and determines the Yoneda embedding
\[
  j \colon \C \to \PSh.
\]

The geometric content of representability is the following standard fact.

\begin{proposition}\label{prop:representable}
For each object $C \in \C$, the projection
\[
  q_C \colon \C^{\op}_{/C} \to \C^{\op}
\]
is a left fibration, and there is a canonical equivalence
\[
  \St_{\C^{\op}}(q_C) \simeq j(C).
\]
Moreover, the object $\id_C \in \C^{\op}_{/C}$ is terminal: for every object
$x \in \C^{\op}_{/C}$, the mapping space
\[
  \Map_{\C^{\op}_{/C}}(x,\id_C)
\]
is contractible.
\end{proposition}

\begin{proof}
The first assertion is the representability statement attached to
straightening--unstraightening for slice left fibrations; equivalently, for each
object $D \in \C$, the fiber of $q_C$ at the corresponding object of $\C^{\op}$ is
canonically equivalent to $\Map_{\C}(D,C)$, compatibly with transport along edges
of $\C^{\op}$. This identifies $\St_{\C^{\op}}(q_C)$ with $j(C)$.

For terminality, let $x \in \C^{\op}_{/C}$. By the definition of the slice
$\infty$-category, the mapping space from $x$ to $\id_C$ is the fiber of the map
from the space of $2$-simplices of $\C$ with terminal vertex $C$ to the space of
pairs of composable edges having composite the edge corresponding to $x$, taken at
the degenerate factorization of that edge through $\id_C$. This fiber is
contractible. Hence $\Map_{\C^{\op}_{/C}}(x,\id_C)$ is contractible.
\end{proof}

\section{Sections and limits}

The Yoneda lemma will be obtained from two formal identifications.

\begin{proposition}\label{prop:sections-pullback}
Let $q \colon X \to K$ and $p \colon Y \to K$ be morphisms in $\mathrm{Cat}_{\infty}$, and
let
\[
\begin{tikzcd}
Y \times_K X \arrow[r] \arrow[d, "p_X"'] & Y \arrow[d, "p"] \\
X \arrow[r, "q"'] & K
\end{tikzcd}
\]
be a pullback square. Then the universal property of the pullback induces a
canonical equivalence of spaces
\[
  \Map_{\mathrm{Cat}_{\infty,/K}}(q,p) \xrightarrow{\simeq} \Gamma(p_X).
\]
If $p$ is a left fibration, then so is $p_X$.
\end{proposition}

\begin{proof}
A point of $\Map_{\mathrm{Cat}_{\infty,/K}}(q,p)$ is, by definition, a morphism
$\alpha \colon X \to Y$ in $\mathrm{Cat}_{\infty}$ together with an equivalence
$p \circ \alpha \simeq q$ in $\Fun(X,K)$. By the universal property of the
pullback, such data are equivalent to a morphism
$\widetilde{\alpha} \colon X \to Y \times_K X$ over $X$, that is, to a section of
$p_X$. This yields the displayed equivalence. Stability of left fibrations under
pullback is standard.
\end{proof}

\begin{proposition}\label{prop:sections-limits}
Let $K$ be an $\infty$-category, and let $r \colon X \to K$ be a left fibration
classified by a functor $G \colon K \to \Sps$. Then there is a canonical
equivalence
\[
  \Gamma(r) \xrightarrow{\simeq} \lim_K G,
\]
functorial in $G$.
\end{proposition}

\begin{proof}
Under straightening, the left fibration $r$ corresponds to $G$. Since
straightening is an equivalence of $\infty$-categories, it identifies the mapping
space
\[
  \Gamma(r)=\Map_{\mathrm{Cat}_{\infty,/K}}(\id_K,r)
\]
with the mapping space in $\Fun(K,\Sps)$ from the terminal object to $G$. By the
universal property of the limit, the latter mapping space is canonically
identified with $\lim_K G$. Functoriality is inherited from functoriality of
straightening and the universal property of the limit functor.
\end{proof}

\section{Finality in the slice}

Fix an object $C \in \C$. Let
\[
  i_C \colon \{\id_C\} \hookrightarrow \C^{\op}_{/C}
\]
denote the inclusion of the full simplicial subset spanned by the terminal object.

\begin{proposition}\label{prop:final}
The functor $i_C$ is final.
\end{proposition}

\begin{proof}
By HTT, \S4.1.3.1, it suffices to show that for each object
$x \in \C^{\op}_{/C}$, the simplicial set
\[
  \{\id_C\} \times_{\C^{\op}_{/C}} (\C^{\op}_{/C})_{x/}
\]
is weakly contractible. By the defining universal property of the overcategory,
this pullback is canonically equivalent to the mapping space
\[
  \Map_{\C^{\op}_{/C}}(x,\id_C).
\]
The latter is contractible by Proposition~\ref{prop:representable}. Hence $i_C$
is final.
\end{proof}

\begin{corollary}\label{cor:restrict-terminal}
For every functor $G \colon \C^{\op}_{/C} \to \Sps$, restriction along $i_C$
induces a canonical equivalence
\[
  \lim_{\C^{\op}_{/C}} G \xrightarrow{\simeq} G(\id_C).
\]
\end{corollary}

\begin{proof}
Since $i_C$ is final, restriction along $i_C$ induces an equivalence
\[
  \lim_{\C^{\op}_{/C}} G \xrightarrow{\simeq} \lim_{\{\id_C\}} G|_{\{\id_C\}}.
\]
The indexing $\infty$-category on the right is contractible with one object, so
its limit is canonically equivalent to the value at that object.
\end{proof}

\section{Yoneda}

The Yoneda lemma is now immediate from the preceding infrastructure.

\begin{theorem}[Yoneda lemma]\label{thm:yoneda}
The bifunctor
\[
  \C^{\op} \times \PSh \to \Sps,
  \qquad
  (C,F) \longmapsto \Map_{\PSh}(j(C),F),
\]
is canonically equivalent to the evaluation bifunctor
\[
  \ev \colon \C^{\op} \times \PSh \to \Sps,
  \qquad
  (C,F) \longmapsto F(C).
\]
Equivalently, for each $C \in \C$ and each $F \in \PSh$ there is a canonical
natural equivalence
\[
  \Phi_{C,F} \colon \Map_{\PSh}(j(C),F) \xrightarrow{\simeq} F(C).
\]
\end{theorem}

\begin{proof}
Fix $(C,F) \in \C^{\op} \times \PSh$. Let
\[
  p \colon \E \to \C^{\op}
\]
be a left fibration classified by $F$, and let
\[
  q_C \colon \C^{\op}_{/C} \to \C^{\op}
\]
be the representable left fibration classified by $j(C)$. Since
$\St_{\C^{\op}}$ is an equivalence, it induces an equivalence on mapping spaces:
\[
  \Map_{\PSh}(j(C),F)
  \simeq
  \Map_{\LFib(\C^{\op})}(q_C,p).
\]
Because $\LFib(\C^{\op})$ is a full subcategory of $\mathrm{Cat}_{\infty,/\C^{\op}}$, the
right-hand side may be computed in the ambient slice $\infty$-category.

Form the pullback square
\[
\begin{tikzcd}
\Ep \arrow[r] \arrow[d, "p^{\prime}"'] & \E \arrow[d, "p"] \\
\C^{\op}_{/C} \arrow[r, "q_C"'] & \C^{\op}.
\end{tikzcd}
\]
By Proposition~\ref{prop:sections-pullback}, there is a canonical equivalence
\[
  \Map_{\LFib(\C^{\op})}(q_C,p) \simeq \Gamma(p^{\prime}).
\]
Let
\[
  F^{\prime} \colon \C^{\op}_{/C} \to \Sps
\]
be the functor classified by $p^{\prime}$. By
Proposition~\ref{prop:sections-limits},
\[
  \Gamma(p^{\prime}) \simeq \lim_{\C^{\op}_{/C}} F^{\prime}.
\]
By Corollary~\ref{cor:restrict-terminal}, restriction along $i_C$ yields a
canonical equivalence
\[
  \lim_{\C^{\op}_{/C}} F^{\prime} \simeq F^{\prime}(\id_C).
\]
Therefore
\[
  \Map_{\PSh}(j(C),F) \simeq F^{\prime}(\id_C).
\]

It remains to identify $F^{\prime}(\id_C)$ with $F(C)$. Since $F^{\prime}$ is the
straightening of $p^{\prime}$, the universal property of straightening identifies
$F^{\prime}(\id_C)$ with the fiber $\Ep_{\id_C}$. On the other hand, by the
universal property of the pullback square defining $\Ep$, the fiber of
$p^{\prime}$ at $\id_C$ is canonically equivalent to the fiber of $p$ at the image
of $\id_C$ under $q_C$, namely the object $C$ of $\C^{\op}$. Thus there are
canonical equivalences
\[
  F^{\prime}(\id_C) \simeq \Ep_{\id_C} \simeq \E_C.
\]
Since $p$ is classified by $F$, straightening identifies $\E_C$ with $F(C)$.
Combining the preceding equivalences gives
\[
  \Map_{\PSh}(j(C),F) \simeq F(C).
\]

Each step is functorial in both variables: $C \mapsto q_C$ is functorial by the
functoriality of the representable left fibration, $F \mapsto p$ is functorial by
unstraightening, pullback is functorial in the diagram, the equivalence between
sections and limits is functorial in the classified diagram, and restriction
along the final functor $i_C$ is functorial in the diagram. Hence the composite
construction refines to an equivalence of bifunctors.
\end{proof}

\begin{corollary}\label{cor:fully-faithful}
The Yoneda embedding
\[
  j \colon \C \to \PSh
\]
is fully faithful. Equivalently, for all objects $C,D \in \C$, the canonical map
\[
  \Map_{\C}(C,D) \to \Map_{\PSh}(j(C),j(D))
\]
is an equivalence of spaces.
\end{corollary}

\begin{proof}
Apply Theorem~\ref{thm:yoneda} with $F=j(D)$. One obtains
\[
  \Map_{\PSh}(j(C),j(D)) \simeq j(D)(C) \simeq \Map_{\C}(C,D).
\]
This equivalence is natural in $(C,D)$, hence identifies the mapping-space
bifunctor of $\C$ with that of the essential image of $j$.
\end{proof}

\begin{remark}
The theorem is not a calculation internal to the presheaf $\infty$-category. It
is the assertion that the representable object $j(C)$ is geometrically realized by
the slice left fibration $\C^{\op}_{/C} \to \C^{\op}$, and that the universal
property of the terminal object $\id_C$ computes sections of its pullback against
an arbitrary left fibration. The Yoneda equivalence is therefore the composite of
straightening, pullback stability, finality, and the universal property of the
limit.
\end{remark}

\end{document}
